% !TEX root =./ECE444_Practice_main.tex
%
% L21 practice set (Array Factor Lab) -- reading HPBW, FNBW and sidelobe level
% off a PHASER beam sweep, predicting a reduced aperture before measuring it,
% and reconciling a measured sidelobe against the array-factor value.

\fullwidth{\textbf{Learning Objective 3.2: } Derive and apply the array factor for a
uniform linear array, and relate its beamwidth and sidelobe structure to what a
beam sweep measures.
\\[4pt]\normalfont\small Show your work. A \textbf{Documentation} line at the end
is required (write \textbf{None} if you did not collaborate).}

\begin{questions}

\question \textbf{Reading a sweep.} A lab group sweeps the PHASER with all eight
elements active and a source at boresight. Working at $10.3\ \ghz$ with element
spacing $d = 14\ \text{mm}$, they record the trace below. The sweep step is
$2.8125\mydeg$ and the peak reads $-4.1$ dBFS.

\begin{center}
\begin{tabular}{r|rrrrrrrrrr}
$\theta$ (\mydeg) & $0.0$ & $2.8$ & $5.6$ & $8.4$ & $11.3$ & $14.1$ & $16.9$ & $19.7$ & $22.5$ & $25.3$ \\ \hline
Power (dBFS) & $-4.1$ & $-4.6$ & $-6.2$ & $-9.1$ & $-14.3$ & $-24.1$ & $-22.0$ & $-17.4$ & $-16.6$ & $-18.1$ \\
\end{tabular}
\end{center}

\begin{parts}

	\begin{minipage}[t]{0.58\linewidth}
		\part Using linear interpolation between the two samples that bracket the
		half-power level, find the half-power beamwidth. The trace is symmetric
		about $0\mydeg$.
		\begin{solution}[1.3in]
		Half power is $3$ dB below the peak, so the level is
		$-4.1 - 3.0 = -7.1$ dBFS. That falls between the samples at $5.6\mydeg$
		($-6.2$) and $8.4\mydeg$ ($-9.1$).
		\eq{ \theta_{-3} &= 5.6 + 2.8\lp\frac{7.1-6.2}{9.1-6.2}\rp = 5.6 + 2.8(0.31) = 6.5\mydeg }
		The trace is symmetric, so the beamwidth is twice this.
		\eq{ \theta_{\text{HP}} = 2(6.5) = 12.9\mydeg }
		Theory gives $13.2\mydeg$, so the sweep reads about three tenths of a
		degree narrow.
		\end{solution}
	\end{minipage}\hfill%
	\begin{minipage}[t]{0.37\linewidth}
		\raggedleft \vspace{12pt}
		$\theta_{\text{HP}}$ = \ansbox[1.3in]{$12.9\mydeg$}
	\end{minipage}

	\begin{minipage}[t]{0.58\linewidth}
		\part Find the first-null beamwidth as the sweep reports it, and compare it
		with $\text{FNBW} = 2\arcsin\lp\lambda/Nd\rp$.
		\begin{solution}[1.4in]
		The deepest sample between the main lobe and the sidelobe is $-24.1$ dBFS
		at $14.1\mydeg$, so the sweep reports
		\eq{ \text{FNBW}_{\text{meas}} = 2(14.1) = 28.2\mydeg }
		With $\lambda = 29.1\ \text{mm}$ and $Nd = 8(14) = 112\ \text{mm}$,
		\eq{ \text{FNBW} &= 2\arcsin\lp\frac{29.1}{112}\rp = 2\arcsin(0.2598) \\ &= 2(15.1\mydeg) = 30.1\mydeg }
		The measured value is $1.9\mydeg$ narrow because the true null at
		$15.1\mydeg$ falls between the samples at $14.1\mydeg$ and $16.9\mydeg$.
		The sweep can only report a null at an angle it actually visited.
		\end{solution}
	\end{minipage}\hfill%
	\begin{minipage}[t]{0.37\linewidth}
		\raggedleft \vspace{12pt}
		$\text{FNBW}_{\text{meas}}$ = \ansbox[1.3in]{$28.2\mydeg$}
		\\[10pt]
		$\text{FNBW}_{\text{calc}}$ = \ansbox[1.3in]{$30.1\mydeg$}
	\end{minipage}

	\begin{minipage}[t]{0.58\linewidth}
		\part Find the first sidelobe level relative to the peak, in dBc.
		\begin{solution}[1.1in]
		The tallest sample beyond the null is $-16.6$ dBFS at $22.5\mydeg$.
		Relative to the peak,
		\eq{ \text{SLL} = -16.6 - (-4.1) = -12.5\ \text{dBc} }
		which is within a few tenths of the $-12.8$ dB the array factor gives
		for eight uniform elements.
		\end{solution}
	\end{minipage}\hfill%
	\begin{minipage}[t]{0.37\linewidth}
		\raggedleft \vspace{12pt}
		SLL = \ansbox[1.3in]{$-12.5$ dBc}
	\end{minipage}

	\part No sample in the sweep reads below about $-24$ dBFS, which is $20$ dB
	below the peak, yet the noise floor on this kit is $23$ dB below the peak.
	Explain in one or two sentences why the deepest sample sits $3$ dB above the
	floor rather than on it.
		\begin{solution}[1.3in]
		The receiver reports signal power plus noise power. At the sample nearest
		the null the array factor has fallen to about $-23$ dB, which is the floor
		level itself, so the two contributions are equal and their sum is $3$ dB
		above either one. A measured null depth is therefore a statement about the
		dynamic range of the measurement rather than about the array.
		\end{solution}

\end{parts}

\newpage

\question \textbf{Predict before you measure.} The same array is reconfigured by
setting some element gains to zero. The spacing $d = 14\ \text{mm}$ and the
frequency $10.3\ \ghz$ do not change, so $\lambda = 29.1\ \text{mm}$ throughout.

\begin{parts}

	\begin{minipage}[t]{0.58\linewidth}
		\part Rx1, Rx2, Rx7 and Rx8 are set to zero, leaving Rx3 through Rx6
		active --- four adjacent elements. Predict the half-power beamwidth from
		$\theta_{\text{HP}} \approx 0.886\lambda/(Nd)$.
		\begin{solution}[1.2in]
		Four active elements span $Nd = 4(14) = 56\ \text{mm}$.
		\eq{ \theta_{\text{HP}} \approx \frac{0.886(29.1)}{56} = 0.4604\ \text{rad} = 26.4\mydeg }
		Solving the array factor exactly gives $27.4\mydeg$; the small-angle step
		inside the $0.886$ rule costs about a degree at this beamwidth.
		\end{solution}
	\end{minipage}\hfill%
	\begin{minipage}[t]{0.37\linewidth}
		\raggedleft \vspace{12pt}
		$\theta_{\text{HP}}$ = \ansbox[1.3in]{$\approx 26\text{-}27\mydeg$}
	\end{minipage}

	\begin{minipage}[t]{0.58\linewidth}
		\part Now only Rx3, Rx4 and Rx5 are left on --- three adjacent elements.
		Predict the half-power beamwidth and the first-null beamwidth.
		\begin{solution}[1.4in]
		Three adjacent elements span $Nd = 3(14) = 42\ \text{mm}$.
		\eq{ \theta_{\text{HP}} \approx \frac{0.886(29.1)}{42} = 0.6138\ \text{rad} = 35.2\mydeg }
		\eq{ \text{FNBW} &= 2\arcsin\lp\frac{29.1}{42}\rp = 2\arcsin(0.6929) \\ &= 2(43.8\mydeg) = 87.7\mydeg }
		The beam is wide enough that the $0.886$ rule is starting to under-predict;
		expect the sweep to read closer to $37\mydeg$.
		\end{solution}
	\end{minipage}\hfill%
	\begin{minipage}[t]{0.37\linewidth}
		\raggedleft \vspace{12pt}
		$\theta_{\text{HP}}$ = \ansbox[1.3in]{$\approx 35\mydeg$}
		\\[10pt]
		FNBW = \ansbox[1.3in]{$87.7\mydeg$}
	\end{minipage}

	\begin{minipage}[t]{0.58\linewidth}
		\part By how many dB should the peak of the three-element sweep sit below
		the peak of the full eight-element sweep, if every element gain that is
		still on remains at 100\%?
		\begin{solution}[1.1in]
		At boresight the element signals add in voltage, so the peak field is
		proportional to the number of active elements.
		\eq{ \Delta = 20\mylog{\frac{3}{8}} = 20\mylog{0.375} = -8.5\ \db }
		\end{solution}
	\end{minipage}\hfill%
	\begin{minipage}[t]{0.37\linewidth}
		\raggedleft \vspace{12pt}
		$\Delta$ = \ansbox[1.3in]{$-8.5\ \db$}
	\end{minipage}

	\part A student proposes reaching the same $35\mydeg$ beamwidth by leaving all
	eight elements on and turning every gain down to $37.5\%$. Explain in one or
	two sentences why that does not work.
		\begin{solution}[1.2in]
		Beamwidth is set by the active aperture $Nd$, not by how much gain the
		elements are given. Turning all eight gains down uniformly keeps the
		aperture at $112\ \text{mm}$, so the beam stays at about $13\mydeg$ and
		only the peak level drops. To widen the beam the aperture itself has to
		shrink.
		\end{solution}

\end{parts}

\newpage

\question \textbf{Reconciling a sidelobe.} A second group measures the first
sidelobe of the full eight-element array at $-11.5$ dBc, while the array factor
for $N = 8$ predicts $-12.8$ dB. Their noise floor sits at $-23$ dBc and the
sweep step is $2.8125\mydeg$.

\begin{parts}

	\begin{minipage}[t]{0.58\linewidth}
		\part The floor adds noise power to the sidelobe power the receiver
		reports. Treating the two as uncorrelated, compute the level the sweep
		should report for a true $-12.8$ dB sidelobe sitting on a $-23$ dBc
		floor.
		\begin{solution}[1.4in]
		Convert both to power ratios and add.
		\eq{ p_{\text{sig}} = 10^{-12.8/10} = 0.0525, \qquad p_{\text{n}} = 10^{-23/10} = 0.0050 }
		\eq{ p_{\text{tot}} = 0.0525 + 0.0050 = 0.0575 }
		\eq{ \text{SLL} = 10\mylog{0.0575} = -12.4\ \text{dBc} }
		The floor alone accounts for about $0.4$ dB of lift.
		\end{solution}
	\end{minipage}\hfill%
	\begin{minipage}[t]{0.37\linewidth}
		\raggedleft \vspace{12pt}
		SLL = \ansbox[1.3in]{$-12.4$ dBc}
	\end{minipage}

	\begin{minipage}[t]{0.58\linewidth}
		\part Their peak reading is also affected. The peak is $23$ dB above the
		floor by definition. By how much does the floor lift the peak reading, and
		which way does that push the reported sidelobe level in dBc?
		\begin{solution}[1.3in]
		Normalize the peak signal power to 1.
		\eq{ p_{\text{tot}} &= 1 + 0.0050 = 1.0050 \\ 10\mylog{1.0050} &= +0.02\ \db }
		The peak lifts by only about $0.02$ dB, which is negligible. Because the
		sidelobe lifts by $0.4$ dB and the peak barely moves, the ratio of the
		two moves \emph{up}, making the reported sidelobe look higher than the
		array factor predicts.
		\end{solution}
	\end{minipage}\hfill%
	\begin{minipage}[t]{0.37\linewidth}
		\raggedleft \vspace{12pt}
		peak lift = \ansbox[1.3in]{$+0.02\ \db$}
	\end{minipage}

	\part The floor explains $0.4$ dB of the $1.3$ dB discrepancy. Name two
	other contributions that plausibly make up the rest, and say for each whether
	it would move the reading up or down.
		\begin{solution}[1.5in]
		First, run-to-run noise on the sidelobe sample itself. The sidelobe is only
		about $10$ dB above the floor, so a single sweep reading scatters a few
		tenths of a decibel in either direction. Second, residual amplitude and
		phase errors across the eight elements after calibration. These break the
		uniform distribution slightly and raise the sidelobes, so that contribution
		moves the reading up. Grid straddle is a third contribution, but it moves
		the reading \emph{down} and by very little here, since the sidelobe peak at
		$21.9\mydeg$ is caught by the sample at $22.5\mydeg$ to within
		$0.05$ dB.
		\end{solution}

\end{parts}

\newpage

\question \textbf{The two-element case.} With only Rx4 and Rx5 active the array
has two elements at $d = 14\ \text{mm}$, and at $10.3\ \ghz$ the array factor is
$AF = \cosp{\psi/2}$ with $\psi = kd\sinp{\theta}$.

\begin{parts}

	\begin{minipage}[t]{0.58\linewidth}
		\part Find the half-power beamwidth exactly, without using the
		$0.886\lambda/(Nd)$ rule.
		\begin{solution}[1.4in]
		Set the array factor to $1/\sqrt{2}$.
		\eq{ \cosp{\psi/2} = \frac{1}{\sqrt{2}} \quad\Rightarrow\quad \frac{\psi}{2} = \frac{\pi}{4} \quad\Rightarrow\quad \psi = \frac{\pi}{2} }
		With $\psi = kd\sinp{\theta}$ and $kd = 2\pi d/\lambda$,
		\eq{ \sinp{\theta} = \frac{\pi/2}{2\pi d/\lambda} = \frac{\lambda}{4d} = \frac{29.1}{4(14)} = 0.5196 }
		\eq{ \theta = 31.3\mydeg \quad\Rightarrow\quad \theta_{\text{HP}} = 62.6\mydeg }
		\end{solution}
	\end{minipage}\hfill%
	\begin{minipage}[t]{0.37\linewidth}
		\raggedleft \vspace{12pt}
		$\theta_{\text{HP}}$ = \ansbox[1.3in]{$62.6\mydeg$}
	\end{minipage}

	\begin{minipage}[t]{0.58\linewidth}
		\part Show that the $0.886\lambda/(Nd)$ rule fails here, and state the
		percentage error it makes.
		\begin{solution}[1.2in]
		\eq{ \theta_{\text{HP}} \approx \frac{0.886(29.1)}{28} = 0.9208\ \text{rad} = 52.8\mydeg }
		\eq{ \text{error} = \frac{52.8 - 62.6}{62.6} = -15.7\% }
		The rule assumes the half-power point is at a small enough angle that
		$\sinp{\theta} \approx \theta$. At $31\mydeg$ that is no longer true.
		\end{solution}
	\end{minipage}\hfill%
	\begin{minipage}[t]{0.37\linewidth}
		\raggedleft \vspace{12pt}
		error = \ansbox[1.3in]{$-15.7\%$}
	\end{minipage}

	\begin{minipage}[t]{0.58\linewidth}
		\part Evaluate $\lambda/2d$ and use it to explain why the lab table lists
		the first-null beamwidth of this array as $180\mydeg$.
		\begin{solution}[1.3in]
		The first null of $\cosp{\psi/2}$ needs $\psi = \pi$, which requires
		\eq{ \sinp{\theta} = \frac{\lambda}{2d} = \frac{29.1}{2(14)} = 1.039 }
		No real angle has a sine greater than 1, so the null lies outside visible
		space and no null appears anywhere in the $\pm 90\mydeg$ scan. With no
		null to measure, the beamwidth between first nulls is quoted as the full
		$180\mydeg$ scan range by convention.
		\end{solution}
	\end{minipage}\hfill%
	\begin{minipage}[t]{0.37\linewidth}
		\raggedleft \vspace{12pt}
		$\lambda/2d$ = \ansbox[1.3in]{$1.039$}
	\end{minipage}

	\begin{minipage}[t]{0.58\linewidth}
		\part How far below its own peak does this two-element pattern fall at
		$\theta = 90\mydeg$? Compare that with the $23$ dB of dynamic range the
		sweep has above the noise floor at eight elements, and say what the trace
		will look like near the edges of the scan.
		\begin{solution}[1.5in]
		At $\theta = 90\mydeg$, $\sinp{\theta} = 1$ and
		$\psi = 2\pi(14)/29.1 = 3.023$ rad, so
		\eq{ \vert AF \vert &= \cosp{1.511} = 0.0597 \\ 20\mylog{0.0597} &= -24.5\ \db }
		The two-element peak is itself $12$ dB below the eight-element peak, so
		it clears the floor by only about $11$ dB. The pattern needs
		$24.5$ dB of range to reach endfire and the measurement has $11$ dB,
		so the trace flattens into the noise well before $\pm 90\mydeg$ and the
		outer part of the roll-off is not visible.
		\end{solution}
	\end{minipage}\hfill%
	\begin{minipage}[t]{0.37\linewidth}
		\raggedleft \vspace{12pt}
		drop at $90\mydeg$ = \ansbox[1.3in]{$-24.5\ \db$}
	\end{minipage}

\end{parts}

\newpage

\question \textbf{Design judgment.} Your section is asked to demonstrate a
$20\mydeg$ beam on the PHASER at $10.3\ \ghz$, using only the Element Gains
controls.

\begin{parts}

	\begin{minipage}[t]{0.58\linewidth}
		\part How many adjacent elements must be active to give a half-power
		beamwidth closest to $20\mydeg$? Give the aperture and the resulting
		beamwidth.
		\begin{solution}[1.4in]
		Solve the beamwidth rule for the aperture.
		\eq{ Nd = \frac{0.886\lambda}{\theta_{\text{HP}}} = \frac{0.886(29.1)}{20\mydeg(\pi/180)} = \frac{25.78}{0.3491} = 73.8\ \text{mm} }
		\eq{ N = \frac{73.8}{14} = 5.3 }
		Element counts are integers, so take $N = 5$: $Nd = 70\ \text{mm}$ and
		\eq{ \theta_{\text{HP}} \approx \frac{0.886(29.1)}{70} = 0.3683\ \text{rad} = 21.1\mydeg }
		Six elements would give $17.6\mydeg$, which is further from the target.
		\end{solution}
	\end{minipage}\hfill%
	\begin{minipage}[t]{0.37\linewidth}
		\raggedleft \vspace{12pt}
		$N$ = \ansbox[1.3in]{$5$ elements}
		\\[10pt]
		$\theta_{\text{HP}}$ = \ansbox[1.3in]{$21.1\mydeg$}
	\end{minipage}

	\begin{minipage}[t]{0.58\linewidth}
		\part What does this cost in peak received level compared with the full
		eight-element array, and how much dynamic range is left above the noise
		floor?
		\begin{solution}[1.3in]
		\eq{ \Delta = 20\mylog{\frac{5}{8}} = 20\mylog{0.625} = -4.1\ \db }
		The floor does not move, so the range above it drops from about $23$ dB
		to
		\eq{ 23 - 4.1 = 18.9\ \db }
		The first sidelobe of a five-element uniform array is near $-12$ dB,
		which still clears the floor comfortably.
		\end{solution}
	\end{minipage}\hfill%
	\begin{minipage}[t]{0.37\linewidth}
		\raggedleft \vspace{12pt}
		$\Delta$ = \ansbox[1.3in]{$-4.1\ \db$}
		\\[10pt]
		range = \ansbox[1.3in]{$18.9\ \db$}
	\end{minipage}

	\part A classmate suggests instead applying a Hann taper to all eight
	elements, since that also widens the beam to about $19.5\mydeg$. State one
	advantage and one disadvantage of the taper against switching elements off,
	and say which you would choose for a demonstration whose purpose is to show
	that beamwidth depends on aperture.
		\begin{solution}[1.6in]
		The taper keeps the full $112\ \text{mm}$ aperture and pushes the
		sidelobes below the noise floor, so the trace is cleaner and the main lobe
		is the only thing on the plot. Its disadvantage for this purpose is that
		it widens the beam by changing the shape of the amplitude distribution,
		not by changing the aperture, so it demonstrates the opposite of the
		intended point. It also costs about $4.7$ dB of peak, which is no better
		than the five-element case. For a demonstration about aperture, switch
		elements off: the beamwidth and the aperture then move together and the
		claim is visible in one control.
		\end{solution}

\end{parts}

\vspace{1.5em}
\noindent\textbf{Documentation:} \rule[-2pt]{0.6\linewidth}{0.4pt}

\end{questions}
